
Input to the Constant-Head (CHD) Package is read from the file that has type ``CHD6'' in the Name File.  Any number of CHD Packages can be specified for a single groundwater flow model; however, an error will occur if a CHD Package attempts to make a GWF cell a constant-head cell when that cell has already been designated as a constant-head cell either within the present CHD Package or within another CHD Package.

Constant-Head input can be specified using lists or arrays.  List-based input for the CHD package is the default, and is described here.  Instructions for specifying array-based CHD input are described in the next section.

List-based input offers several advantages over the array-based input for specifying constant heads.  First, multiple list entries can be specified for a single cell.  This makes it possible to divide a cell into multiple areas, and assign a different constant head for each area.  In this case, the user would likely specify an auxiliary variable to serve as a multiplier.  This multiplier would be calculated by the user and provided in the input file as the fractional cell area for the individual constant head entries.  Another advantage to using list-based input for specifying constant heads is that boundnames can be specified.  Boundnames work with the Observations capability and can be used to sum constant heads for entries with the same boundname.  A disadvantage of the list-based input is that one cannot easily assign constant heads to the entire model without specifying a list of model cells.  For this reason \mf also supports array-based input for constant heads.

In previous MODFLOW versions, it was not possible to convert a constant-head cell to an active cell.  Once a cell was designated as a constant-head cell, it remained a constant-head cell until the end of the end of the simulation.  In \mf a constant-head cell will become active again if it is not included as a constant-head cell in subsequent stress periods.

Previous MODFLOW versions allowed specification of SHEAD and EHEAD, which were the starting and ending prescribed heads for a stress period.  Linear interpolation was used to calculate a value for each time step.  In \mf only a single head value can be specified for any constant-head cell in any stress period.  The time-series functionality must be used in order to interpolate values to individual time steps.  

\vspace{5mm}
\subsubsection{Structure of Blocks}
\vspace{5mm}

\noindent \textit{FOR EACH SIMULATION}
\lstinputlisting[style=blockdefinition]{./mf6ivar/tex/gwf-chd-options.dat}
\lstinputlisting[style=blockdefinition]{./mf6ivar/tex/gwf-chd-dimensions.dat}
\vspace{5mm}
\noindent \textit{FOR ANY STRESS PERIOD}
\lstinputlisting[style=blockdefinition]{./mf6ivar/tex/gwf-chd-period.dat}
\packageperioddescription

\vspace{5mm}
\subsubsection{Explanation of Variables}
\begin{description}
\input{./mf6ivar/tex/gwf-chd-desc.tex}
\end{description}

\vspace{5mm}
\subsubsection{Example Input File}
\lstinputlisting[style=inputfile]{./mf6ivar/examples/gwf-chd-example.dat}

\vspace{5mm}
\subsubsection{Available observation types}
CHD Package observations are limited to the simulated constant head flow rate (\texttt{chd}). The data required for the CHD Package observation type is defined in table~\ref{table:gwf-chdobstype}. Negative and positive values for an observation represent a loss from and gain to the GWF model, respectively.

\begin{longtable}{p{2cm} p{2.75cm} p{2cm} p{1.25cm} p{7cm}}
\caption{Available CHD Package observation types} \tabularnewline

\hline
\hline
\textbf{Model} & \textbf{Observation type} & \textbf{ID} & \textbf{ID2} & \textbf{Description} \\
\hline
\endhead

\hline
\endfoot

\input{../Common/gwf-chdobs.tex}
\label{table:gwf-chdobstype}
\end{longtable}

\vspace{5mm}
\subsubsection{Example Observation Input File}
\lstinputlisting[style=inputfile]{./mf6ivar/examples/gwf-chd-example-obs.dat}

